\section{Research Plan}
\label{sec:plan}

We organize the empirical program as a partially-observed sequential decision problem over experiments: each phase is selected to maximize value-of-information per unit of compute, so that cheap, high-diagnostic measurements gate the expensive ones. The ordering is deliberate, and it is now disciplined by a result we obtained while writing this plan. The thesis of the paper is that the optimization gain $\gain$ of training-free RL---Gibbs tilting $\qs \propto \qo\exp(\Rwd/\beta)$ of a frozen base $\qo$---is governed by the \emph{reward range} $\spread = \max_{\operatorname{supp}\qo}\Rwd - \min_{\operatorname{supp}\qo}\Rwd$ of the (action-space, reward) pair (OSA-1, \Cref{thm:hoeffding}); to leading order in $1/\beta$ the operative discriminator is the base-measure variance $\operatorname{Var}_{\qo}(\Rwd)/(2\beta)$ (a remark refines the certified range-based ceiling to this variance form). This is a property of the \emph{reward structure}, not of model class: a high-range verifiable reward yields large gain \emph{even on a single frozen model} (our own single-model Operator-B results, e.g.\ SLURP intent $+0.330$, \Cref{tab:opB}), while a flat reward---instruction-conditioning of a frozen embedding read-out---yields vanishing gain regardless of how the model is wrapped. Context isolation makes gain \emph{additive} across blocks (OSA-2, \Cref{thm:gainprod}) and each non-degenerate block contributes \emph{strictly positive} gain (\texttt{gain\_pos\_of\_nonconstant}: reward non-constant $\Rightarrow \gain>0$, sorry-free in Lean), but \Cref{thm:qstarprod} (\texttt{qstar\_product}) proves the context-isolated optimum equals the monolithic single-model optimum on the product reward $\Rwd=\sum_i\Rwd_i$---so \emph{isolation per se adds nothing}; the only new gain comes from introducing additional \emph{non-degenerate verifiable rewards}. In particular, our high-range content/intent gains are \emph{single-model} optima, \emph{not} evidence that agentic decomposition recovers anything a single model could not. Whether a frozen omni embedding read-out affords such non-degenerate per-block rewards for our flagship paralinguistic factors is therefore an empirical question (\Cref{conj:spreadfloor}), and the answer we have so far is humbling.

\paragraph{The result that re-scoped this plan.} We ran the paired-delta bootstrap that \Cref{sec:feasibility} flagged as outstanding Phase-0 work, now committed at \texttt{projects/speech-mllm-omni-embedding-rl/\_repro/emotion\_pool\_paired\_v2.json}. With the read-out layer selected on a dev split (no oracle test-layer peeking) and a paired-delta bootstrap over $5$ seeds ($42,7,123,2024,31337$), the emotion-pooling Operator-A effect is a \emph{null} at the across-seed level: the mean is $\Delta=+0.037$, but the $95\%$ across-seed $t$-confidence interval is $[-0.043,+0.116]$ and spans $0$, so the effect is not significant.\footnote{As a diagnostic only, the per-seed paired CI excludes $0$ in $2/5$ seeds; this seed-vote count is \emph{not} the inferential statistic---the across-seed CI is, and it does not exclude $0$.} The previously reported single-seed $+0.097$ was the best seed under \emph{oracle test-layer selection}---a winner's-curse / selection-leak artifact, not a stable effect. In parallel, a frozen linear/kNN probe on the omni \emph{content} bi-encoder (Omni-Embed) sits near chance ($\approx 0.04$) for speaker identity at \emph{every} one of its $37$ layers. We are careful about the scope of this negative: it is established on \emph{one content bi-encoder} via a linear/kNN read-out probe, and it bounds only that Operator-A read-out path; it does \emph{not} bound the proposed \emph{generative} Operator-B policy, which we have not yet tested on paralinguistics. The honest conclusion is that the spread-floor conjecture holds for content/intent, while for the flagship paralinguistic factors the one frozen embedding we have probed exposes no non-degenerate per-block reward (emotion null, speaker at chance). We therefore re-scope: high-range content/intent verifiable rewards yield large \emph{single-model} gains, whereas whether \emph{agentic decomposition} adds anything for speaker/affect is unresolved and---by \Cref{thm:qstarprod}---can be settled only by introducing genuinely new non-degenerate verifiable rewards, which remain untested. This re-scoping is load-bearing for every baseline and admission band below; the plan is written so that a null paralinguistic result is a reportable outcome, not a failure of the instrument.

We first close the loop on a single GPU with frozen weights (Phase~0), because no downstream claim can be trusted until the verifiable-reward gate is demonstrably faithful end-to-end and the single-model spread is measured honestly. We then build the missing measurement instrument (Phase~1), since the central question---whether reward-guided inference-time optimization over $\cZ$ can \emph{compound} competence across sessions---cannot be answered without a benchmark that does not yet exist, and that benchmark must be \emph{falsifiable against a real, teethed baseline}. Only then do we run the minimal self-evolving loop (Phase~2) and ablate its controls (Phase~3). Throughout, the only quantities that drive selection are verifiable rewards $\Rwd$, and the locked test set is touched exactly once. We stress at the outset that the benchmark, and the baseline the agent must beat, are the load-bearing risks in this plan: a poorly constructed cross-session benchmark, or a strawman baseline, can manufacture an apparent compounding effect out of artifacts, so a substantial part of Phase~1 is devoted to ruling out such artifacts and to pre-registering a classic-pipeline baseline that the agent must \emph{exceed} before any agentic claim is made.

\subsection{Phase 0: Closed-Loop Bring-Up}
\label{sub:phase0}

The objective of Phase~0 is a verified, reproducible inference-time loop on the in-place RTX~5090 (Blackwell, sm\_120; torch \texttt{cu128}) with \emph{no weight updates and no structural change}. The frozen stack couples an omni-embedding model (Omni-Embed-Nemotron, \citealp{omniembed2025,omniembedaudio2026}) for the embedding operator $\cZA$ with Qwen3-Omni \citep{qwen3omni} served via GGUF/\texttt{llama.cpp} for the generative operator $\cZB$. Concretely, this realizes the two-operator action space $\cZ = \cZA \cup \cZB$ used throughout: $\cZA$ selects pooling/probe read-outs of a frozen representation; $\cZB$ is reward-guided decoding (best-of-$N$, reranking, soft-BoN) over text. In light of the result above, we make a structural commitment up front: the omni embedding is reserved for \emph{content} read-out and retrieval, where its spread is real, while \emph{speaker} and \emph{affect} memory keys are carried by classic speech encoders (ECAPA-TDNN speaker embeddings \citealp{ecapatdnn2020} and a dedicated SER classifier), not by the omni embedding. The ``two-omni'' system is thus an omni \emph{policy} (Operator-B) plus an omni \emph{content} memory, paired with classic paralinguistic encoders---we do \emph{not} claim the omni embedding performs load-bearing paralinguistic retrieval, because on our own content-bi-encoder probe it is at chance on speaker.

We reproduce two known frozen baselines as correctness anchors: SER accuracy $\approx 0.49$ and intent accuracy $\approx 0.25$ under zero-shot prompting of the backbone. Matching these within bootstrap tolerance certifies that data loading, prompt templating (\texttt{speechrl\_common.models}), and decoding are wired correctly. Because these anchors come from public corpora that may overlap a backbone's pretraining data, we treat them as wiring checks only, not as evidence of generalization, and we subject every public-dataset preliminary number to the contamination analysis described below before quoting it. We then verify the verifiable-reward gate: for ASR we instrument WER-based exact rewards, and for SER/SID we instrument label-exact rewards (\texttt{speechrl\_common.rl}), confirming that (i) best-of-$N$ selection under $\Rwd$ strictly dominates greedy on a held-in probe, consistent with the regret/winrate scaling of soft and hard best-of-$N$ \citep{verdun2025softbon,beirami2024bon,aminian2025smoothing}, and (ii) the implied search distribution $\qz$ tilts the base $\qo$ toward the tilted optimum $\qs \propto \qo \exp(\Rwd/\beta)$ of the objective $\Fobj{q} = \Exp_q[\Rwd] - \beta\,\KL{q}{\qo}$ as $\beta$ is lowered, with the partition function $\Zpart$ and objective $\Fobj{\cdot}$ tracked as in the KL-as-Bayes view \citep{korbak2022kl}. Critically, Phase~0 also \emph{re-confirms the spread map} that drives the rest of the plan: we log per-utterance $\gain$ and $\spread$ for each candidate reward and report which rewards are high-range (content/intent best-of-$N$) and which are flat (paralinguistic instruction-conditioning), using the dev-selected, multi-seed paired-delta protocol that produced the null $+0.037$ emotion result above---never oracle test-layer selection. The exit criterion is a green end-to-end run logging $\gain$ and $\spread$ per utterance to local MLflow, a passing \texttt{pytest} smoke suite, and a committed, dev-selected paired-delta artifact for every paralinguistic readout claim. Phase~0 buys the highest information per dollar: it eliminates instrumentation confounds, and it forces every paralinguistic gain through the same selection-leak-proof, across-seed-CI protocol before any benchmark or loop claim is made.

\paragraph{Contamination analysis of preliminary public-dataset numbers.} Before any of the Phase~0 anchor numbers (or other figures computed on existing public datasets) enter the paper, each is passed through a contamination screen, not only the newly constructed benchmark. The screen combines (i) train-cutoff reasoning, comparing each corpus's release date against the backbone's reported pretraining cutoff and flagging any corpus plausibly in-distribution; (ii) $n$-gram / transcript overlap between corpus references and any recoverable pretraining-adjacent text, reported as an overlap rate per split; and (iii) a canary / memorization probe that checks whether the backbone can complete held-out references above chance without the audio \citep{contaminationsurvey2024,taskcontamination2024}. Numbers that fail any screen are reported as wiring checks with an explicit contamination caveat rather than as generalization evidence. We make no claim that this screen is exhaustive; it is a documented best effort whose results we report alongside the affected numbers.

\subsection{Phase 1: A Speech Cross-Session Paralinguistic Memory Benchmark}
\label{sub:phase1}

\paragraph{The gap.} Existing long-horizon memory benchmarks are text-only (\citealp{longmemeval2024,locomo2024}; the agentic memory line \citealp{amber2026,anatomymem2026}), and audio memory work targets within-document or RAG retrieval \citep{wavrag2025,voxrag2025,audiomc2025}. No raw-audio, \emph{cross-session}, paralinguistically-keyed memory benchmark exists. This is the intended contribution: a benchmark whose queries can only be answered by retaining \emph{who spoke} and \emph{how they sounded} across sessions, not by transcribing content. Whether a cheaply constructed benchmark can faithfully isolate this competence is itself an open question, and---given the negative paralinguistic evidence above---we treat the construction below as a hypothesis to be validated rather than a settled instrument, and we equip it with a real baseline that the agent must beat before any agentic claim is asserted.

\paragraph{Primary baseline: the strongest simple pipeline the agent must exceed.} The single most important addition to this plan is a \emph{teethed} primary baseline with a \emph{pre-registered, locked} accumulation rule, and it must be the genuinely strongest \emph{simple} pipeline---not a strawman---or the comparison cannot isolate agentic gain. We split the baseline \emph{by factor}, because the identity (trait) and affect-change (state) sub-problems have different strongest-simple solutions.

\emph{Identity (trait) baseline: ECAPA $+$ PLDA with full multi-enrollment.} For speaker identity we pre-register, in the locked release artifact, a frozen \textbf{ECAPA-TDNN} front-end \citep{ecapatdnn2020,xvector2021} scored by \textbf{PLDA} (equivalently, cosine with adaptive symmetric normalization, AS-Norm) under \emph{full multi-enrollment}: all of a speaker's enrolled utterances across sessions are combined under the PLDA model, not merely averaged into a recent-$k$ centroid. This is a deliberate \emph{upgrade} from the cosine running-centroid dictionary we first considered, and the reason is structural: the candidate agentic mechanism of Phase~2---Bayesian multi-session evidence integration, i.e.\ accumulating a per-speaker posterior over identity from per-session ECAPA read-outs treated as noisy observations---\emph{is essentially PLDA}. PLDA is exactly the generative model that treats each embedding as a latent speaker factor plus channel/session noise and integrates evidence across enrollments. A cosine-centroid baseline would therefore let the agent book a ``gain'' that is nothing more than PLDA-over-cosine (proper channel-aware multi-enrollment scoring)---a classical, non-agentic improvement---so it \emph{cannot isolate agentic gain from PLDA-over-cosine}. Only by pre-registering ECAPA$+$PLDA-with-multi-enrollment as the baseline do we force the agentic store to demonstrate that it adds something \emph{beyond} correct Bayesian evidence integration---concretely: graph-structured curation of which enrollments to keep, \textsc{update}/\textsc{delete} decay of stale keys, a trust region on the per-episode store change, and explicit change handling (speaker re-enrollment on detected drift). The agentic store must \emph{strictly exceed} ECAPA$+$PLDA multi-enrollment, or its identity claim is unsupported.

\emph{Affect-change (state) baseline: per-session-independent SER.} For affect we pre-register a \textbf{per-session-independent SER} baseline---each session's affect is predicted from that session alone by the frozen SER labeler, with \emph{no} cross-session state, optionally smoothed by \textbf{majority vote with exponential decay} for a running-state variant---together with an explicit \textbf{change-detection threshold} on affect drift. The agentic change-point store must beat this at \emph{detecting affect changes} across sessions: the test is whether a curated cross-session store detects a genuine affect change (or its absence) better than independently labeling each session and differencing. As with identity, the baseline's exact \textsc{update}/\textsc{delete} and change-detection-threshold rule is locked in the release artifact.

We require the self-evolving agent to \emph{strictly exceed} the relevant baseline on each factor, with every contrast tested by a paired bootstrap on per-item deltas \citep{bisani2004bootstrap} over $\geq 5$ seeds with across-seed CIs and FWER correction across the sweep. We pre-register two \emph{separate} falsifiers so that a null on one does \emph{not} over-generalize to ``the agentic claim for paralinguistics'' as a whole: \textbf{(trait/identity)} if multi-session curation does not exceed ECAPA$+$PLDA multi-enrollment with a paired-bootstrap CI excluding $0$, the agentic claim \emph{on identity} is recorded as unsupported; \textbf{(state/affect-change)} if the change-point store does not beat the per-session-independent SER baseline at detecting affect changes with a paired-bootstrap CI excluding $0$, the agentic claim \emph{on affect change} is recorded as unsupported. A null on identity leaves the affect-change question open, and vice versa. These baselines are deliberately strong precisely because the one frozen content bi-encoder we probed is at chance on speaker: the classic pipeline encodes exactly the paralinguistic competence the omni \emph{content} embedding lacks, so beating it---\emph{beyond} mere PLDA evidence integration---is the real test of whether the \emph{agentic} composition (memory $\times$ policy $\times$ verifiable gate) adds anything. Both baselines, with their locked rules, ship in the release artifact so any reader can re-run the comparison.

\paragraph{Minimal spec.} (i) \emph{Stimuli}: multi-session spoken dialogues delivered as raw audio, with $\geq 3$ sessions per persona separated by simulated time gaps. (ii) \emph{Queries}: speaker-ID-keyed and SER-keyed (e.g., ``which speaker was angry in the previous call,'' ``has this caller's affect changed since session~1''), so that the intended answer is a function of paralinguistic state, not lexical content---guarding against the read-not-listen shortcut \citep{voxparadox2026,listenortranscribe2025}. (iii) \emph{Ground truth}: SID labels are treated as verifiable; SER labels are \emph{not} treated as ground truth from a frozen labeler (see the SER gating paragraph), and only the human-validated, high-confidence SER subset is admitted for exact rewards or reporting. (iv) \emph{Probes}: a knowledge-update probe (a speaker's emotion/role changes across sessions, testing \textsc{update}/\textsc{delete} semantics \citealp{mem0_2025}), a contamination probe \citep{contaminationsurvey2024,taskcontamination2024}, and a memory-poisoning probe \citep{agentpoison2024,mempoison2026}.

\paragraph{Query-generalization split (recall is not memory).} A positive cross-session slope is meaningless if it is answer-cache recall. We therefore split queries into a \emph{repeated} set (the same question re-asked across sessions) and a \emph{novel-query} set (new questions about the \emph{same} speaker's stored paralinguistic state, never asked before). Stored state must \emph{transfer} to the novel queries; we report the compounding slope \emph{separately} on repeated vs.\ novel queries, and treat the novel-query slope as the primary compounding evidence. A slope that appears only on repeated queries is reported as caching, not memory.

\paragraph{Leakage control: raw embeddings only, no gold labels in store.} To make a positive slope un-attributable to key-value recall, the memory store is forbidden from writing gold labels: only \emph{raw embeddings} (ECAPA speaker vectors, SER logits/embeddings, omni content vectors) may be persisted, never the benchmark's gold SID/SER answer. Thus a positive slope cannot be produced by reading back a stored label; it must arise from re-reading the acoustic state. The store schema is validated automatically to contain no gold-label field.

\paragraph{Construction and its realism risks.} The cheap construction we initially considered---re-keying an existing diarized corpus (SLURP/SUPERB-style intent+speaker material, \citealp{slurp2020,superb2021,dynsuperb2023}) into ``multi-session dialogues''---has a serious validity flaw we make explicit: stitching single-speaker clips of the same speaker into separate ``sessions'' yields \emph{channel-matched, exact-audio} continuity, so cross-session speaker linking can be solved by low-level acoustic/channel fingerprinting rather than by any genuine paralinguistic memory, trivializing the competence the benchmark is meant to measure. We therefore require \emph{genuine} cross-session variation as an admission condition: either (a) the same persona is re-recorded under different recording conditions (device, room, codec, background), or (b) we source a real repeated-caller corpus in which a speaker recurs naturally across genuinely distinct sessions. Channel-matched exact-audio re-keying is used, if at all, only as an explicitly labeled \emph{upper-bound/ceiling} control, never as the headline cross-session set.

\paragraph{Numeric admission bands (DER / SER / SV-EER), fixed in advance.} To prevent ``disqualify if near-trivial'' from being post-hoc tunable, we pre-register numeric bands and freeze them before construction. The speaker-linking sub-problem is admitted only if the frozen ECAPA speaker-verification \emph{equal-error rate} (SV-EER) on the constructed cross-session pairs lies in $[8\%,30\%]$: an EER below $5\%$ signals channel fingerprinting (trivial; disqualified), above $40\%$ signals no recoverable speaker signal (degenerate; disqualified). Diarization difficulty is admitted only if the frozen-diarizer \emph{DER} lies in $[10\%,30\%]$ (below $5\%$ trivial, above $40\%$ unusable). The SER sub-problem is admitted only if frozen-labeler-vs-human balanced accuracy on a held-out calibration slice lies in $[0.35,0.75]$---above the $k$-way chance floor but well below ceiling, so the task is neither trivial nor saturated. These three bands are reported on every constructed split; any split outside a band is rejected or relabeled a ceiling control. To quantify residual leakage we additionally report DER and SER error on the constructed set so the difficulty of the speaker-linking and affect sub-problems is measured rather than assumed.

\paragraph{Listen-vs-read control as a primary metric, run on the agent's own transcripts.} Every paralinguistic query must pass a transcript-only / listen-only control before admission: a strong text LLM given only the gold transcripts (no audio) must perform at chance on the paralinguistic queries; queries solvable from transcript alone are removed. We harden this control in three ways. First, we run the transcript condition on the agent's \emph{own ASR transcripts} in two variants---a \emph{surface-marked} version preserving the cues a careless transcriber might leak (capitalization, ``!!!'', \texttt{[laughter]}, disfluency markers) \emph{and} a \emph{normalized} version that strips them---so a gap that vanishes only under normalization is exposed as surface leakage, not listening. Second, we \emph{balance topic across speakers}, so that lexical topic cannot proxy for speaker or affect identity. Third, we add a \emph{positive control}: a model independently known to attend to acoustics is run on the same queries, and we require it to show a substantial audio$-$transcript gap; if even the positive control cannot, the gap reflects benchmark difficulty rather than a single model's deafness, and we say so. We report the \emph{audio-vs-transcript gap}---audio-conditioned minus transcript-only performance---as a \emph{primary} metric: a benchmark with a small or zero gap is, by construction, not measuring listening. Construction uses \emph{frozen} diarization (pyannote, \citealp{pyannotecode,personalvad2019,xvector2021}) and a frozen SER labeler only as \emph{proposal} tools whose outputs are audited and, where used for ground truth, human-validated; their raw outputs are never the final labels.

\paragraph{SER-keyed ground truth, spontaneous and acted affect reported separately.} We do not treat a frozen SER labeler as ground truth. Any SER-keyed write into memory, and any SER-keyed reward, is gated behind an SER-confidence threshold: low-confidence affect predictions are abstained on rather than committed. For the \emph{evaluated} subset, SER labels are obtained by human annotation, with multiple annotators per item; we report inter-annotator agreement (Krippendorff's $\alpha$/Fleiss' $\kappa$) and exclude items below a pre-registered agreement floor. Acted affect is known to overstate SER reliability, so we source at least one \emph{spontaneous} (non-acted) affect corpus---in-the-wild, naturally elicited speech---and report it \emph{separately} from acted CREMA-D, never pooled. For the knowledge-update probe specifically we report the \emph{fraction of affect-change items that survive the IAA floor}: if few affect-change items reach inter-annotator agreement, the update probe is under-powered and we report it as such rather than quoting a slope on a handful of items. Exact SER rewards are computed only against human-agreed, high-confidence labels; the frozen labeler serves solely to pre-filter and scale annotation, never to define correctness.

\paragraph{Evaluation design.} Scoring uses verifiable rewards only, reported with paired bootstrap confidence intervals over the per-utterance deltas \citep{bisani2004bootstrap}, standard frozen splits \citep{gorman2019splits}, and an audited memory-eval rubric \citep{anatomymem2026}. We report, as primary and \emph{per factor}, the agent-minus-baseline paired-bootstrap contrast (agent vs.\ ECAPA$+$PLDA for identity, agent vs.\ per-session-independent SER for affect change); as secondary, DER/SV-EER/SER admission bands, the audio-vs-transcript gap, the repeated-vs-novel-query slopes, the spontaneous-vs-acted affect split, \emph{maintenance latency} (wall-clock cost of the memory \textsc{add}/\textsc{update} cycle), and we run every metric across $\geq 2$ backbones (Qwen3-Omni and Qwen2.5-Omni, \citealp{qwen25omni}) to separate benchmark signal from backbone idiosyncrasy. Phase~1 is sequenced second because it is the precondition for any compounding claim, yet is independent of the loop's internals, so it can be built and frozen---and its realism controls, numeric bands, and baselines passed---before Phase~2 risks overfitting to it.

\subsection{Phase 2: Minimal Self-Evolving Loop}
\label{sub:phase2}

Phase~2 instantiates the smallest loop that could plausibly compound: a paralinguistic \emph{memory} store (raw-embedding keys with \textsc{add}/\textsc{update}/\textsc{delete} \citealp{mem0_2025,memoryos2025}, ECAPA/SER keys for paralinguistics and omni keys for content, never gold labels), a \emph{skill gate} that admits a candidate skill only when a held-out verifiable reward improves \citep{skillsbench2026,coevoskills2026}, and a \emph{trust region} that bounds the per-episode change in the search distribution via $\KL{\qz}{\qo} \le \epsilon$, capping the inference-time tilt to limit reward over-optimization \citep{gao2023overopt,skalse2022rewardhacking}. The loop is gradient-free: credit is assigned by reward comparison across rollouts \citep{jitrl2026,ahmadian2024rloo}, and selection over $\cZB$ uses reranking/MBR among $N$ samples \citep{ichihara2025mbr,mbrasr2025}. SER-keyed writes inherit the confidence gating and abstention discipline of Phase~1, so unreliable affect predictions do not silently become persistent state. The loop's design is consistent with \Cref{thm:qstarprod}: because the context-isolated optimum equals the monolithic single-model optimum on the product reward, the loop adds value only insofar as it introduces \emph{additional non-degenerate verifiable rewards}---which, by our Phase-0 evidence, are abundant for content/intent and scarce for the paralinguistic read-out we probed. Concretely, Phase~2 tests candidate paralinguistic mechanisms \emph{by factor}: for identity, \textbf{Bayesian multi-session evidence integration} (treating each session's ECAPA read-out as a noisy observation and accumulating a per-speaker posterior over identity) against the pre-registered ECAPA$+$PLDA multi-enrollment baseline; for affect, a \textbf{change-point store} over per-session SER read-outs against the per-session-independent SER baseline. Because the identity mechanism is essentially PLDA (\Cref{sub:phase1}), only the ECAPA$+$PLDA baseline can isolate genuine agentic gain from PLDA-over-cosine---the agentic store must add graph curation / decay / trust region / change handling \emph{on top of} correct evidence integration. A null on the identity test disqualifies the \emph{agentic claim for identity}, and a null on the affect-change test disqualifies the \emph{agentic claim for affect change}, but \emph{neither} disqualifies the other factor nor the underlying theory (\Cref{thm:qstarprod} already predicts no gain absent a genuinely new non-degenerate reward). We therefore expect and pre-register an \emph{asymmetric}, \emph{per-factor} outcome.

We \emph{pre-register} success and kill criteria before running, separately by factor. \textbf{Success (content/intent):} a positive, bootstrap-significant cross-episode slope on intent/spoken-QA reward, on the \emph{novel-query} split, that \emph{exceeds the classic-pipeline baseline} under matched per-speaker evidence with a paired-bootstrap CI excluding $0$, with the gate's admit-rate stable and $\KL{\qz}{\qo}$ inside the trust region. \textbf{Trait/identity hypothesis (pre-registered as a null):} we do \emph{not} predict that multi-session curation exceeds ECAPA$+$PLDA multi-enrollment on SID; a null or baseline-not-exceeded result on the novel-query SID split is an \emph{expected and reportable} outcome, given the negative read-out evidence and \Cref{conj:spreadfloor}. \textbf{State/affect-change hypothesis (pre-registered as a null):} we likewise do \emph{not} predict that the change-point store beats the per-session-independent SER baseline at detecting affect changes on the human-validated SER subset; a null there is \emph{separately} expected and reportable, and does \emph{not} generalize to the identity claim. \textbf{Kill:} no slope after a fixed episode budget, admit-rate collapse, trust-region saturation (drift), or---decisively---failure of a factor's agentic mechanism to exceed \emph{its} pre-registered baseline (ECAPA$+$PLDA for identity, per-session-independent SER for affect change), which we record as \emph{not supporting the agentic claim} on \emph{that factor only}. The headline measurement is \emph{compounding}: reward as a function of episode index. The \emph{primary} contrast is agent-minus-baseline where \emph{both receive the identical accumulated per-speaker evidence}---the same enrolled utterances, in the same order---so the comparison isolates the agentic update rule from the sheer amount of evidence. We further impose an \textbf{evidence-matched / fixed-audio-budget control}: the total per-speaker enrolled audio (equivalently, the number of enrolled utterances) is held \emph{constant} across episode index, so a positive compounding slope cannot be driven by mere acoustic-data accumulation rather than by self-evolution of the memory. The vs-empty-memory single-shot contrast (best-of-$N$ with non-persistent memory) is \emph{demoted to a secondary diagnostic}: it cannot separate self-evolution from accumulation and so cannot, on its own, establish compounding. The compounding measurement isolates cross-session state from raw test-time compute \citep{snell2024ttscompute,zuo2025ttrl} and from a trivial accumulating dictionary. Because a spurious slope could arise from benchmark artifacts or answer caching, this measurement is interpretable only on a Phase~1 set that has passed the DER/SV-EER/SER admission bands, the audio-vs-transcript-gap control, the raw-embeddings-only leakage control, and the novel-query split.

\subsection{Phase 3: Scale and Ablate}
\label{sub:phase3}

Phase~3 stress-tests each control by removal, with pre-registered directional predictions. (i) \emph{Remove the skill gate}: we predict SkillsBench-style degradation as unvetted skills accumulate \citep{skillsbench2026,ace2025}. (ii) \emph{Remove the trust region}: we predict drift, i.e., $\KL{\qz}{\qo}$ growth correlated with falling locked-test reward, the inference-time analogue of reward hacking \citep{gao2023overopt,remembering2026}. (iii) \emph{Append-only vs.\ curated memory}: we predict append-only error propagation and context collapse \citep{expfollowing2025,ace2025}, with curated \textsc{update}/\textsc{delete} recovering it \citep{mem0_2025}. (iv) \emph{Omni vs.\ classic paralinguistic keys}: we predict that swapping the omni embedding in for the ECAPA/SER keys \emph{degrades} speaker/affect compounding toward chance, directly testing the negative read-out finding that motivated this plan. Each ablation is a single-variable contrast against the Phase~2 configuration, scaled across $N\in\{4,8,16,32\}$ and across backbones to chart the reward-vs-compute and reward-vs-tilt frontiers. These predictions are hypotheses; a null or reversed result is reported as such rather than reinterpreted post hoc.

\subsection{Evaluation Protocol}
\label{sub:protocol}

The protocol is fixed in advance and applies uniformly. (1) \emph{Rewards}: verifiable only (WER/exact-match/label-exact), with SER rewards restricted to the human-validated, confidence-gated subset; no LLM-judge in any selection or reporting path, avoiding self-preference and position bias \citep{selfpreference2024,positionbias2024}. (2) \emph{Test discipline}: model selection and all tuning---including read-out layer/pooling selection---occur on a disjoint dev split only \citep{dwork2015reusable,roelofs2019overfitting}; \emph{oracle test-layer selection is prohibited}, having already produced the spurious single-seed $+0.097$ emotion gain that the dev-selected re-run corrected to a null $+0.037$ (across-seed CI $[-0.043,+0.116]$, spanning $0$). After dev-only selection, the locked test set is evaluated in a single, untouched pass at the end of each phase. (3) \emph{Demonstration hygiene}: in-context demonstrations are drawn disjoint from both dev and test, since demonstrations alter format more than accuracy \citep{min2022rethinking,alice2026}, and we report true few-shot results \citep{perez2021truefewshot}. (4) \emph{Uncertainty}: every comparison is run over $\geq 5$ seeds with across-seed confidence intervals \citep{henderson2018drlmatters}, and deltas between conditions are tested with a \emph{paired} bootstrap on the per-item differences \citep{bisani2004bootstrap,dror2018hitchhiker}---reporting, as the emotion result does, the across-seed mean and CI as the inferential statistic (with the fraction of seeds whose paired CI excludes $0$ given only as a diagnostic, never as the headline)---with a power analysis fixing the minimum detectable effect \citep{card2020power}. (5) \emph{Multiplicity}: family-wise error-rate (FWER) correction is applied across the config/ablation sweep \emph{and across the two pre-registered per-factor baseline contrasts}, so that selecting among many configurations (or many read-out layers) does not inflate the apparent significance of the single test-pass deltas. (6) \emph{Contamination}: the Phase~0 contamination screen (train-cutoff reasoning, $n$-gram/transcript overlap, canary) is applied to \emph{every} public-dataset number, not only the constructed benchmark, and its results are reported alongside the affected figures. (7) \emph{Negative results are reported as first-class outcomes}---the null emotion result and the at-chance speaker probe (on one content bi-encoder) are reported in the paper, not suppressed, and per-factor nulls are reported per factor without over-generalization---and the full configuration follows the reproducibility checklist \citep{pineau2020reproducibility,dodge2019showyourwork}.

\subsection{Milestones, Compute, and Reproducibility}
\label{sub:milestones}

Compute is already provisioned: an in-place RTX~5090 with torch \texttt{cu128}, all five models (including the Qwen3-Omni GGUF and Omni-Embed-Nemotron \citealp{omniembedweights}) and all 28 datasets resident on WSL ext4 disk, frozen to a pinned manifest. Every run is reproducible by construction: pinned data revisions $+$ fixed seeds $+$ a single \texttt{reproduce:} command logged to MLflow, so any reported number maps to one invocation---including the committed paired-delta artifact \texttt{emotion\_pool\_paired\_v2.json} that re-scoped this plan.

\begin{table}[t]
\centering
\small
\begin{tabular}{llll}
\hline
\textbf{Milestone} & \textbf{Phase} & \textbf{Deliverable} & \textbf{Exit gate} \\
\hline
M0 & 0 & Closed loop on frozen $\cZA{+}\cZB$ & SER$\,{\approx}\,0.49$, intent$\,{\approx}\,0.25$ reproduced (wiring) \\
M1 & 0 & Reward-gate $+$ contamination screen $+$ spread map & best-of-$N$ $>$ greedy, CI-sig.; dev-selected paired-delta logged \\
M2 & 1 & Cross-session benchmark $+$ ECAPA$+$PLDA / SER baselines & genuine variation; DER/SV-EER/SER in-band; locked rules shipped \\
M3 & 1 & Listen-vs-read $+$ human-SER $+$ novel-query split & audio$-$transcript gap $>0$ (incl.\ positive control); IAA floor met \\
M4 & 2 & Minimal self-evolving loop & pre-registered success/kill adjudicated per factor \\
M5 & 2 & Compounding vs.\ baselines (evidence-matched) & slope exceeds per-factor baseline (ECAPA$+$PLDA identity; per-session SER affect) under matched evidence (paired CI) \emph{or} documented null \\
M6 & 3 & Control ablations (incl.\ omni-vs-classic keys) & gate/TR/append-only/key-type contrasts complete \\
M7 & 3 & Frontier sweep ($N$, backbone) & reward-vs-compute/tilt curves $+$ release \\
\hline
\end{tabular}
\caption{Value-of-information--ordered milestones. Each gate is a verifiable, bootstrap-tested criterion; the locked test set is touched once per phase, after dev-only selection (no oracle test-layer selection), with FWER correction across the sweep and the two per-factor contrasts. The compounding gate (M5) requires \emph{exceeding the per-factor baselines---ECAPA$+$PLDA multi-enrollment for identity, per-session-independent SER for affect change---under matched per-speaker evidence}, not merely a positive slope, and a null is adjudicated per factor.}
\label{tab:milestones}
\end{table}

The release artifact bundles the benchmark builder (including the realism controls: cross-session variation requirement, numeric DER/SV-EER/SER admission bands, listen-vs-read filtering on own ASR transcripts with surface-marked and normalized variants, the spontaneous-affect split, and the SER human-annotation pipeline), the two per-factor baselines with their pre-registered locked rules---the ECAPA$+$PLDA identity baseline (frozen ECAPA-TDNN front-end, PLDA/AS-Norm scoring under full multi-enrollment) and the per-session-independent SER affect baseline (SER majority-with-decay, explicit change-detection threshold)---the raw-embeddings-only memory schema, the loop configuration, pinned manifests, the contamination-screen scripts, the committed paired-delta artifacts, and per-milestone \texttt{reproduce:} commands, so that all of \Cref{tab:milestones} can be re-derived end-to-end on equivalent hardware. This ordering guarantees that the expensive compounding study (M4--M7) is only attempted after the instrument and its baselines (M2--M3) and the gate and spread map (M0--M1) are independently certified---including their artifact, contamination, and selection-leak controls---maximizing the information returned for the fixed single-GPU budget, and ensuring that an honest negative on either the identity or the affect-change agentic claim is a publishable, per-factor result of this plan rather than a hidden failure.